\chapterimage{figs/chapterhead.png}
\chapter{Results, Traces, and Reports}
\label{chap:results-traces-and-reports}

This chapter documents the result path that already exists in the repository.
It does not invent a new analysis workflow. Instead, it explains how runtime
traces, report text, aggregated statistics, and grouped plots are produced by
the current GUI and kernel code.

\noindent\textbf{Objective.} Explain how users inspect traces, simulation
reports, aggregated statistics, and run-to-run comparisons after a simulation.

\noindent\textbf{Audience.} Users who need to read simulation output, compare
runs, and understand what the reported numbers do and do not prove.

\noindent\textbf{Scope.} Cover the trace console, the simulation and reports
panes, the aggregated results table, the grouped result plots, the replication
reporting flow, and the current limits of interpretation and export.

\noindent\textbf{Status.} Partially validated by code inspection. The chapter
matches the current widget names, reporter calls, and statistics columns, but
the figures are still structural placeholders until the final screenshots are
produced.

The chapter follows this sequence:
Figure~\ref{fig:results-trace},
Figure~\ref{fig:results-report},
Figure~\ref{fig:results-table},
Figure~\ref{fig:results-plots},
Figure~\ref{fig:results-replication-comparison}, and
Figure~\ref{fig:results-data-flow}.

\section{Trace surfaces and report routing}
\label{sec:results-trace-surfaces}

The GUI currently separates three textual destinations for runtime output. The
trace controller routes generic traces and fatal or recoverable errors to
\texttt{textEdit\_Console}, simulation-event traces to
\texttt{textEdit\_Simulation}, and report traces to
\texttt{textEdit\_Reports}. That split is implemented in
\texttt{TraceConsoleController} and wired from the main window.

The report path is explicit in the kernel. The default simulation reporter
emits replication and full-simulation summaries. Those report methods write
lines with \texttt{TraceManager::Level::L2\_results}, so the GUI can display
them as report output instead of ordinary console chatter.

The simulation pane is for live execution feedback. It shows event-level text,
including the lines emitted while a replication is running or advancing through
events. The reports pane is for structured summary text, not for free-form
debugging.

\begin{figure}[htbp]
\centering
% FIGURE-SPEC-BEGIN
% ID: fig:results-trace
% Source of truth:
%   - source/applications/gui/genesys/controllers/TraceConsoleController.cpp
%   - source/applications/gui/genesys/mainwindow.ui
%   - source/applications/gui/genesys/mainwindow.h
%   - source/applications/gui/genesys/mainwindow_simulator.cpp
% Purpose:
%   Show the live trace surfaces used during simulation execution.
% Required visual content:
%   - the main console trace pane
%   - the simulation trace pane
%   - a representative live event or results line
% Accessibility description:
%   A screenshot that distinguishes console traces from simulation traces during a run.
% Validation:
%   The visible pane labels and routing must match the current GUI code.
% FIGURE-SPEC-END
\figureplaceholder{figs/generated/results-trace}
\caption{Live trace surfaces used while a simulation is running.}
\label{fig:results-trace}
\end{figure}

As shown in Figure~\ref{fig:results-trace}, the runtime trace surfaces are
separated by purpose. That separation matters because a trace line that explains
a simulation event should not be read as a final statistical result.

\section{Replication and simulation reports}
\label{sec:results-reports}

The current reporter emits two related textual summaries:
\begin{itemize}
\item a replication report, produced after each replication when reporting is
enabled;
\item a simulation report, produced after the full run finishes.
\end{itemize}

The replication report starts and ends with explicit markers that identify the
replication number. Inside that block, the reporter groups statistics and
counters by parent type and parent name, then prints the columns currently used
by the kernel: name, element count, minimum, maximum, average, variance,
standard deviation, variation coefficient, confidence interval half-width, and
confidence level. Counters are printed as counts only.

The simulation report reuses the same grouping logic across the complete run.
It shows the same statistic columns, but the numbers are aggregated across the
full simulation rather than a single replication. The current code can also show
simulation controls and, when enabled, a list of simulation responses.

The chapter stays conservative about the word \textit{export}. In the current
GUI, the primary artifact is the rendered report text in the reports pane. For
offline result-file analysis, the repository also provides
\texttt{SimulationResultsDatasetParser} in \texttt{source/tools}, which reads
text datasets and Genesys \texttt{Record} files for analysis outside the main
simulation window.

\begin{figure}[htbp]
\centering
% FIGURE-SPEC-BEGIN
% ID: fig:results-report
% Source of truth:
%   - source/kernel/simulator/SimulationReporterDefaultImpl1.cpp
%   - source/kernel/simulator/SimulationReporterDefaultImpl1.h
%   - source/applications/gui/genesys/controllers/TraceConsoleController.cpp
%   - source/applications/gui/genesys/mainwindow.ui
%   - source/applications/gui/genesys/mainwindow_simulator.cpp
% Purpose:
%   Show a real replication or full-simulation report rendered in the reports pane.
% Required visual content:
%   - the reports text pane
%   - begin/end report markers
%   - grouped statistics or counters
% Accessibility description:
%   A report screenshot that makes the structured summary output readable.
% Validation:
%   The headings and row structure must match the reporter implementation.
% FIGURE-SPEC-END
\figureplaceholder{figs/generated/results-report}
\caption{Structured replication or full-simulation report in the reports pane.}
\label{fig:results-report}
\end{figure}

As illustrated in Figure~\ref{fig:results-report}, the report text is not a
free-form log. It is a structured summary, so the surrounding text should be
read as statistical reporting rather than event narration.

\section{Aggregated results table}
\label{sec:results-table}

The aggregated table in the reports area is populated from the current
simulation aggregate list returned by the kernel. The GUI configures the table
with these columns:
\begin{itemize}
\item Type;
\item ParentType;
\item ParentName;
\item Name;
\item NumElements;
\item Min;
\item Max;
\item Average;
\item Variance;
\item StdDev;
\item VarCoef;
\item HalfWidthCI;
\item ConfidenceLevel.
\end{itemize}

The table is filled from the current simulation's aggregate data definitions.
When a row corresponds to a \texttt{StatisticsCollector}, the GUI reads the
statistics object and prints the full descriptive set. When a row corresponds to
a \texttt{Counter}, the GUI prints the counter value and leaves the descriptive
statistics cells blank. Parent type and parent name are derived from the owning
model element so the table can be read in the same hierarchy used by the report
text.

\begin{figure}[htbp]
\centering
% FIGURE-SPEC-BEGIN
% ID: fig:results-table
% Source of truth:
%   - source/applications/gui/genesys/mainwindow.cpp
%   - source/kernel/simulator/model/ModelSimulation.cpp
%   - source/kernel/simulator/model/ModelSimulation.h
%   - source/kernel/simulator/essentialPlugins/StatisticsCollector.h
%   - source/kernel/simulator/essentialPlugins/Counter.h
% Purpose:
%   Show the aggregated results table populated after a simulation run.
% Required visual content:
%   - the reports results table
%   - statistic rows and/or counter rows
%   - the current column set
% Accessibility description:
%   A screenshot that exposes the summary metrics in tabular form.
% Validation:
%   The visible columns must match the table header setup in code.
% FIGURE-SPEC-END
\figureplaceholder{figs/generated/results-table}
\caption{Aggregated simulation results table.}
\label{fig:results-table}
\end{figure}

Figure~\ref{fig:results-table} is the quickest way to inspect the numeric
summary for a run. It should be read together with the report text, because the
table makes the values compact but does not explain the model context by itself.

\section{Grouped result plots}
\label{sec:results-plots}

The plots area is built dynamically inside the \texttt{tabReportsPlots} tab.
The GUI wraps the content in a scroll area and creates one chart per result
category. The charting widget is a custom \texttt{ResultsBoxPlotWidget}, not a
generic plotting component from an external tool.

The chart grouping is conservative. Items are grouped by parent type. Each item
is labeled with either the result name itself or a \textit{parent.name} form
when a parent exists. The chart draws the minimum and maximum as whiskers and
the average as a central horizontal line. The filled box is a presentation aid;
it should not be read as a strict quartile box plot unless the underlying code
changes to compute quartiles explicitly.

The result plots are intentionally summary plots. They are useful for spotting
wide ranges, missing stability, or obviously skewed responses, but they do not
replace a statistical test or a design-of-experiments analysis.

\begin{figure}[htbp]
\centering
% FIGURE-SPEC-BEGIN
% ID: fig:results-plots
% Source of truth:
%   - source/applications/gui/genesys/mainwindow.cpp
%   - source/applications/gui/genesys/mainwindow.ui
%   - source/kernel/simulator/model/ModelSimulation.cpp
% Purpose:
%   Show the grouped result charts rendered in the reports plots tab.
% Required visual content:
%   - the plots tab or scroll area
%   - one or more grouped charts
%   - visible labels for the plotted result categories
% Accessibility description:
%   A chart screenshot that makes the summary trend by category visible.
% Validation:
%   The chart type and grouping must match the current custom plot widget.
% FIGURE-SPEC-END
\figureplaceholder{figs/generated/results-plots}
\caption{Grouped result charts rendered from simulation aggregates.}
\label{fig:results-plots}
\end{figure}

As illustrated in Figure~\ref{fig:results-plots}, the plot view complements the
table view. The table is exact and compact; the plot is visual and comparative.

\section{Comparing replications}
\label{sec:results-replication-comparison}

Replication comparison in the current code is primarily textual and summary
based. The replication report makes each replication boundary explicit, which
lets the user compare the same statistic across consecutive replications when
the model and settings are held constant.

If the user needs a file-based comparison, the repository also contains the
dataset parser used by the analysis tools. That parser understands Genesys
\texttt{Record} files and other text datasets with replication metadata, so the
results can be inspected outside the simulation run itself.

The important limit is methodological: a replication-by-replication comparison
is only meaningful when the runs are comparable. Different seeds, different
input data, or different model settings change the interpretation. The chapter
therefore treats replication comparison as a descriptive aid, not as proof of
equivalence.

\begin{figure}[htbp]
\centering
% FIGURE-SPEC-BEGIN
% ID: fig:results-replication-comparison
% Source of truth:
%   - source/kernel/simulator/SimulationReporterDefaultImpl1.cpp
%   - source/tools/Statistics/SimulationResultsDatasetParser.cpp
%   - source/tools/README_tools.md
%   - source/applications/gui/genesys/mainwindow.cpp
% Purpose:
%   Show a comparison between two or more replications using the current reporting output.
% Required visual content:
%   - at least two replication summaries or equivalent comparison rows
%   - a visible shared metric across replications
%   - the replication identifiers or headings
% Accessibility description:
%   A comparison screenshot that makes run-to-run differences readable.
% Validation:
%   The compared numbers must come from the current replication report or dataset parser.
% FIGURE-SPEC-END
\figureplaceholder{figs/generated/results-replication-comparison}
\caption{Comparison of results across replications.}
\label{fig:results-replication-comparison}
\end{figure}

Figure~\ref{fig:results-replication-comparison} should be read as a comparison
aid only. It does not claim that the reported mean or interval alone proves a
global model property.

\section{Data to report pipeline and limits}
\label{sec:results-data-flow}

The current data path is:
\begin{enumerate}
\item the simulation kernel updates statistics, counters, and responses;
\item the simulation reporter renders replication or full-simulation summaries;
\item the trace controller routes report text into the reports pane;
\item the GUI fills the results table and grouped plots from aggregate data;
\item the analysis tools can load text datasets and \texttt{Record} files for
offline inspection.
\end{enumerate}

The pipeline is complete enough to read a real run, but it is not a promise that
every metric is automatically exported in a user-friendly file format. The
chapter therefore distinguishes three layers:
\begin{itemize}
\item live trace output;
\item structured report text and aggregate tables;
\item offline dataset loading for analysis.
\end{itemize}

The current limits are equally important:
\begin{itemize}
\item the plots are summary visuals, not a statistical inference engine;
\item the report text is generated from the current kernel aggregates, not from
an external audited data warehouse;
\item the chapter does not claim a dedicated one-click export path that is absent
from the current GUI.
\end{itemize}

\begin{figure}[htbp]
\centering
% FIGURE-SPEC-BEGIN
% ID: fig:results-data-flow
% Source of truth:
%   - source/kernel/simulator/SimulationReporterDefaultImpl1.cpp
%   - source/applications/gui/genesys/controllers/TraceConsoleController.cpp
%   - source/applications/gui/genesys/mainwindow.cpp
%   - source/applications/gui/genesys/mainwindow_simulator.cpp
%   - source/tools/Statistics/SimulationResultsDatasetParser.cpp
% Purpose:
%   Show the flow from simulation data to report text, table, plots, and offline analysis.
% Required visual content:
%   - kernel statistics and counters
%   - reporter output
%   - GUI text/table/plot sinks
%   - offline dataset analysis entry point
% Accessibility description:
%   A flow diagram that explains how numeric results reach the user-visible panes.
% Validation:
%   Every arrow and sink must correspond to an actual code path or documented tool.
% FIGURE-SPEC-END
\figureplaceholder{figs/generated/results-data-flow}
\caption{Data flow from simulation aggregates to GUI and offline analysis.}
\label{fig:results-data-flow}
\end{figure}

As shown in Figure~\ref{fig:results-data-flow}, the reporting path is layered.
That layering is useful because it keeps live runtime feedback separate from the
aggregated numbers used for reading and comparison.
